%-*-latex-*-
%-*-latex-*-
\newcommand\COURSE{ciss350}
\newcommand\ASSESSMENT{q10707}
\newcommand\ASSESSMENTTYPE{Quiz}
\newcommand\POINTS{\textwhite{xxx/xxx}}

\input{myquizpreamble}
\input{yliow}
\input{\COURSE}
%-*-latex-*-

%-*-latex-*-

%-*-latex-*-
\renewcommand\TITLE{\ASSESSMENTTYPE \ \ASSESSMENT}

\renewcommand\EMAIL{}
\renewcommand\AUTHOR{eagraham1@cougars.ccis.edu}

\textwidth=6in
\begin{document}
\topmatterthree

Here's the \verb!BTNode! class:
\begin{Verbatim}[frame=single,fontsize=\footnotesize]
class BSTNode
{
public:
    int key_;
    BTNode * parent_;
    BTNode * left_, * right_;
};
\end{Verbatim}

For the following questions,
here's an example on how to describe a binary tree in text.
For this tree
\begin{console}[fontsize=\footnotesize]
            (10)
          /      \
      ( 4)        (25)
     /           /    \
 ( 1)        (12)      (27)
                           \
                            (29)
\end{console}
you write
\begin{console}[fontsize=\footnotesize]
[10, [4, 25]]
[4, [1, None]]
[25, [12, 27]]
[1, [None, None]]
[12, [None, None]]
[27, [None, 29]]
[29, [None, None]]
\end{console}

%------------------------------------------------------------------------------
\nextq
In moving a pointer \verb!p! up one node at a time during
balancing, you need to keep track on the pointer that is \lq\lq behind"
\verb!p! and you also need to remember the path backward by two steps, whether
the 2-step path downward is a
LEFT-LEFT or LEFT-RIGHT or RIGHT-LEFT or RIGHT-RIGHT.
For instance if \verb!p! is here:
\begin{console}[fontsize=\footnotesize]
            (10)
          /      \
      ( 4)        (25)
     /           /    \
 ( 1)    p-> (12)      (27)
                           \
                            (29)
\end{console}
and if it climbs up by one step, you have
\begin{console}[fontsize=\footnotesize]
            (10)
          /      \
      ( 4)        (25) <- p, down = LEFT, downdown = NONE
     /           /    \
 ( 1)     q-> (12)      (27)
                           \
                            (29)
\end{console}
and if it climbs one step up, you have
\begin{console}[fontsize=\footnotesize]
            (10) <- p, down = RIGHT, downdown = LEFT
          /      \
      ( 4)        (25) <- q
     /           /    \
 ( 1)        (12)      (27)
                           \
                            (29)

\end{console}
The constants \verb!LEFT!, \verb!RIGHT!, \verb!NONE! are defined in the code
fragment below.
If \verb!p! is the root node, after calling \verb!climbup!,
you do the obvious -- for instance
\verb!p! becomes \verb!NULL!, \verb!q! points to the root node,
\verb!down! is \verb!NONE!, and \verb!downdown! is the previous value of
\verb!down!.
\\
\ANSWER
\begin{answercode}
const int NONE = -1;
const int LEFT = 0;
const int RIGHT = 1;
void climbup(BTNode *& p, BTNode *& q, int & down, int & downdown)
{
  BTNode * oldp = p;

  p= p->parent_;

  q = p;

  down = downdown;

  if(p == NULL)
  {
    down = NONE;
  }
  else if(p->left_ = q)
  {
    down = LEFT;
  }
  else
  {
    down = RIGHT;
  }
}
\end{answercode}

%------------------------------------------------------------------------------
\nextq
For this tree:
\begin{console}[fontsize=\footnotesize]
            (10)
          /      \
      ( 4)        (25)
     /           /    \
 ( 1)        (12)      (27)
                           \
                            (29)
\end{console}
What is the tree after performing a right rotation at the node
with key value of $10$?
\\
\ANSWER
\begin{answercode}
[4, [1, 10]]
[1, [None, None]]
[10, [None, 25]]
[25, [12, 27]]
[12, [None, None]]
[27, [None, 29]]
[29, [None, None]]
\end{answercode}

%------------------------------------------------------------------------------
\nextq
For this BST:
\begin{console}[fontsize=\footnotesize]
            (10)
          /      \
      ( 4)        (25)
     /           /    \
 ( 1)        (12)      (27)
                           \
                            (29)
\end{console}
insert 30 and then balance it (if necessary).
Write the resulting tree in text below.
\ANSWER
\begin{answercode}
[25, [10, 29]]
[10, [4, 12]]
[29, [27, 30]]
[4, [1, None]]
[12, [None, None]]
[27, [None, None]]
[30, [None, None]]
[1, [None, None]]
\end{answercode}

%------------------------------------------------------------------------------
\nextq
For this BST:
\begin{console}[fontsize=\footnotesize]
            (10)
          /      \
      ( 4)        (25)
     /           /    \
 ( 1)        (12)      (27)
                           \
                            (29)
\end{console}
insert 2 and then balance it (if necessary).
Write the resulting tree in text below.
\ANSWER
\begin{answercode}
[10, [4, 25]]
[2, [1, 4]]
[25, [12, 27]]
[1, [None, None]]
[4, [None, None]]
[12, [None, None]]
[27, [None, 29]]
[29, [None, None]]
\end{answercode}

%------------------------------------------------------------------------------
%-*-latex-*-
\newpage
\input{instructions}
\end{document}

